\chapter{Couche 5 — L'infrastructure}

Jusqu'ici, tout suppose des serveurs déjà installés. Cinq briques pour que ces
serveurs se créent et se configurent tout seuls, à l'identique, autant de fois
qu'on veut.

\begin{figure}[htbp]
\centering
\resizebox{0.92\textwidth}{!}{%
\begin{tikzpicture}[font=\small,
  b/.style={draw=grisligne,fill=grisfond,rounded corners=3pt,minimum height=11mm,
            minimum width=3.2cm,align=center},
  f/.style={-{Stealth[length=2.5mm]},gristexte,thick}]
\node[b] (r) at (0,0) {\textbf{Rien}\\\textcolor{gristexte}{\scriptsize un dossier}};
\node[b] (v) at (4.2,0) {\textbf{Vagrant}\\\textcolor{gristexte}{\scriptsize 4 VM nues}};
\node[b] (a) at (8.4,0) {\textbf{Ansible}\\\textcolor{gristexte}{\scriptsize 4 VM configurées}};
\node[b] (p) at (12.6,0) {\textbf{La plateforme}\\\textcolor{gristexte}{\scriptsize qui tourne}};
\draw[f] (r) -- node[above,font=\scriptsize]{\cmd{vagrant up}} (v);
\draw[f] (v) -- node[above,font=\scriptsize]{\cmd{ansible-playbook}} (a);
\draw[f] (a) -- node[above,font=\scriptsize]{démarrage} (p);
\end{tikzpicture}}
\caption{Deux commandes séparent un dossier vide d'une plateforme en
fonctionnement. Tout ce qui se passe entre les deux est décrit dans des fichiers
versionnés.}
\end{figure}

% -----------------------------------------------------------------------------
\outil{Vagrant}{brique 29 · créer les machines virtuelles}

\quoi
Un outil qui décrit des machines virtuelles dans un fichier texte et les crée en
une commande. \cmd{vagrant up} les monte, \cmd{vagrant destroy} les efface.

\pourquoi
L'alternative est de créer les VM à la main dans l'interface de VirtualBox ou de
VMware : choisir l'image, la mémoire, le réseau, installer le système. Une demi-
heure par machine, et personne ne saura ensuite ce qui a été coché.

Terraform aurait pu faire le travail, mais il vise les fournisseurs de cloud
publics ; pour des VM locales, Vagrant est plus direct et c'est l'outil que le
sujet nomme.

\role
Le \fichier{Vagrantfile} décrit un contrôleur en 192.168.56.10 et trois workers
en .11, .12, .13, avec leur mémoire et leurs cœurs. Le nombre de workers et les
tailles sont réglables sans toucher au fichier :

\begin{lstlisting}[language=bash]
WORKERS=2 CTRL_MEM=1536 WORKER_MEM=1280 vagrant up
\end{lstlisting}

\begin{retenir}
Rendre les tailles paramétrables coûte trois lignes et change tout pour qui
reprend le projet : on peut le faire tourner sur un portable modeste. Un fichier
qui n'offre qu'une seule taille finit toujours par être modifié à la main, donc
désynchronisé du dépôt.
\end{retenir}

\apport
Le parc entier se crée en une commande, et se détruit aussi vite. On peut tout
casser pour vérifier que tout se remonte.

% -----------------------------------------------------------------------------
\outil{Les providers : VMware, VirtualBox, Parallels, QEMU}{brique 30 · exécuter les VM}

\quoi
Le \emph{provider} est le logiciel de virtualisation que Vagrant pilote. Vagrant
ne virtualise rien lui-même : il donne des ordres à l'un d'eux.

\pourquoi
Un seul provider aurait suffi... sur une seule architecture. Le projet devait
tourner sur des postes Intel et sur des Mac à puce Apple.

\textbf{VirtualBox} est le provider par défaut et le plus répandu, mais son
support ARM reste expérimental : sur un Mac récent, il ne convient pas.
\textbf{VMware Fusion} fonctionne nativement sur Apple Silicon.
\textbf{Parallels} aussi, mais il est payant. \textbf{QEMU} est libre et
fonctionne partout, au prix de limitations.

Déclarer les quatre laisse le même fichier fonctionner sur n'importe quel poste.

\role
Chaque provider a son bloc dans le \fichier{Vagrantfile}, avec ses réglages
propres.

\begin{piege}{Chaque provider a ses limites}
\cmd{vagrant-qemu} ne sait pas créer de réseau privé : l'accès SSH doit être
redirigé sur \cmd{localhost}. Le fichier le détecte et ajuste la configuration.
Sans cela, le déploiement réussirait et rien ne serait joignable --- la pire
sorte de panne, celle qui ne dit pas son nom.
\end{piege}

\apport
Le projet se monte sur le poste de n'importe quel membre du groupe, et sur
l'ordinateur de la salle de soutenance.

% -----------------------------------------------------------------------------
\outil{Ansible}{brique 31 · configurer les machines}

\quoi
Un outil qui se connecte en SSH aux machines et applique un \textbf{état décrit
en YAML} : tel paquet installé, tel fichier présent avec tels droits, tel
service démarré.

\pourquoi
L'alternative est un script shell. La différence n'est pas cosmétique.

Un script décrit des \textbf{actions} : installe, crée, copie, redémarre. Le
relancer refait tout, écrase, redémarre inutilement, et échoue si quelque chose
existe déjà.

Ansible décrit un \textbf{état voulu}. Il regarde, compare, et n'agit que si
nécessaire. C'est la propriété d'\textbf{idempotence} : appliquer deux fois
donne le même résultat qu'appliquer une fois.

Les concurrents --- Puppet, Chef, Salt --- exigent un agent installé sur chaque
machine. Ansible n'a besoin que de SSH et de Python, déjà présents. C'est aussi
l'outil que le sujet nomme.

\begin{notion}{Idempotent ne veut pas dire « ne fait rien »}
Cela veut dire « vérifie l'état, et n'agit que s'il diffère de celui qui est
voulu ». La nuance compte : c'est elle qui permet de relancer le déploiement
sans crainte, à n'importe quel moment, y compris sur un parc à moitié
configuré.

La preuve se lit dans le compte rendu : au second passage, \cmd{changed=0} sur
tous les nœuds.
\end{notion}

\role
Neuf rôles, quatre playbooks. Chaque rôle est autonome, avec ses variables par
défaut, ses fichiers et ses gabarits.

\begin{center}
\small
\begin{tabularx}{\textwidth}{@{}lX@{}}
\toprule
\textbf{Rôle} & \textbf{Ce qu'il garantit} \\
\midrule
\cmd{security\_hardening} & le pare-feu UFW, avec les ports du groupe \\
\cmd{controller\_key} & le compte de service et sa clé SSH de pilotage \\
\cmd{docker} & Docker Engine sur les workers, le client seul sur le contrôleur \\
\cmd{tls\_cert} & un certificat propre à chaque nœud \\
\cmd{worker} & l'accès SSH restreint du contrôleur \\
\cmd{docker\_images} & les images des machines louées, construites sur place \\
\cmd{worker\_agent} & l'agent d'inscription, en service systemd \\
\cmd{reverse\_proxy} & nginx en HTTPS devant Gunicorn \\
\cmd{webapp} & l'application, le Faucheur, le watchdog, les secrets \\
\bottomrule
\end{tabularx}
\end{center}

\begin{piege}{Un rôle appelé avec des paramètres différents s'exécute deux fois}
Ansible dédoublonne les rôles... sauf si on les appelle avec des paramètres
différents : il les considère alors comme distincts. Le rôle \cmd{docker} était
appelé directement \emph{et} tiré en dépendance par \cmd{docker\_images}, avec
un paramètre différent --- il s'exécutait deux fois, installant puis
désinstallant le démon.

La correction : passer par une \textbf{variable de play}, qui s'applique aux
deux appels et les fait fusionner. Nous l'avons vérifié par l'expérience, en
instrumentant le rôle pour qu'il journalise chaque exécution.
\end{piege}

\apport
Quatre machines nues deviennent une plateforme en fonctionnement, en une
commande, à l'identique. Et le déploiement peut être rejoué sans risque.

% -----------------------------------------------------------------------------
\outil{Ansible Vault}{brique 32 · les secrets dans Git}

\quoi
Un fichier chiffré en AES-256 qu'Ansible déchiffre au moment de jouer le
playbook, avec une phrase de passe fournie à l'exécution.

\pourquoi
Les secrets --- clé de session, clé du coffre, mot de passe de la base, jeton des
agents --- doivent arriver sur le serveur. Trois mauvaises façons de s'y
prendre :

\textbf{Les écrire en clair dans le dépôt} : \cmd{gitleaks} les trouverait, et
n'importe qui ayant accès au dépôt aussi.

\textbf{Les passer en variables d'environnement à la main} : il faut s'en
souvenir, et ils finissent dans l'historique du shell.

\textbf{Les mettre dans un gestionnaire externe} --- HashiCorp Vault, AWS Secrets
Manager : la bonne solution en entreprise, mais c'est un service de plus à
déployer et à maintenir.

Ansible Vault est le compromis juste pour ce projet : un fichier,
\textbf{versionnable tel quel}, illisible sans la phrase de passe.

\role
\fichier{3\_ansible/group\_vars/insacloud/vault.yml} contient trois secrets. Le
rôle \cmd{webapp} vérifie même leur forme avant de déployer --- la clé Fernet
doit faire exactement 44 caractères --- plutôt que de laisser l'application
échouer au démarrage.

La phrase de passe vit dans un fichier \cmd{.vault\_pass} en \cmd{0600}, exclu
du dépôt par le \fichier{.gitignore}.

\apport
Les secrets voyagent avec le code sans jamais être lisibles. Et le déploiement
reste automatisable : pas d'étape manuelle de copie de clés.

% -----------------------------------------------------------------------------
\outil{L'inventaire dynamique}{brique 33 · découvrir les machines}

\quoi
Un script exécutable qui, au lieu de lire une liste figée de machines, interroge
Vagrant et produit l'inventaire en JSON.

\pourquoi
L'inventaire statique est un fichier qui liste les machines et leurs adresses.
Il a deux défauts : il faut le tenir à jour quand le parc change, et il code en
dur le chemin de la clé SSH --- qui dépend du provider.

C'est un piège réel : l'inventaire du projet contenait
\cmd{vagrant\_provider=virtualbox}. En passant à VMware, Ansible cherchait la
clé au mauvais endroit.

\role
\fichier{3\_ansible/inventaire\_vagrant.py} interroge \cmd{vagrant ssh-config} et
reprend le chemin de clé que Vagrant \emph{déclare lui-même}. Il ne découvre que
les machines réellement démarrées.

\begin{lstlisting}[language=bash]
ansible-playbook -i inventaire_vagrant.py site.yml --ask-vault-pass
\end{lstlisting}

\apport
Le même dépôt se déploie sur les quatre providers sans modification, et sur un
parc de taille variable. C'est le quatrième ajout au-delà du sujet.

% -----------------------------------------------------------------------------
\section*{Ce que la plateforme sait faire après la couche 5}

Elle se construit entièrement toute seule : quatre machines virtuelles, leur
système, leurs paquets, leurs certificats, leurs services, ses secrets. Deux
commandes depuis un dossier vide.

Il reste à garantir qu'elle ne se dégrade pas.
